\documentclass[10pt,a4paper]{amsart}
\usepackage[margin=19mm]{geometry}
\usepackage{amsmath,amssymb,mathtools}
\usepackage[scr=rsfs,cal=euler]{mathalfa}
\usepackage{xfrac}
\usepackage[unicode,hidelinks,bookmarks=false]{hyperref}
\newcommand{\ie}{i.e.\ }
\newcommand{\cf}{cf.\ }
\newcommand{\resp}{respectively\ }
\newcommand{\Subsection}{\subsection}
\providecommand{\nobreakdash}{\nobreak\hbox{-}}
\setlength{\emergencystretch}{2em}
\multlinegap0pt
\usepackage{amssymb}
\newtheorem{lemma}{Lemma}
\newcommand{\R}{\mathbb{R}}
\newcommand{\Z}{\mathbb{Z}}
\title{Sums of squares on curves and surfaces}
\author{Bartłomiej Bychawski, Bartosz Głowacki, Tomasz Kowalczyk}
\date{Source: arXiv:2606.22401v1 (CC BY 4.0)}
\begin{document}
\maketitle

\noindent\emph{Source and licence.} Sums of squares on curves and surfaces. Version arXiv:2606.22401v1 (CC BY 4.0), distributed by arXiv under the Creative Commons Attribution 4.0 International licence, http://creativecommons.org/licenses/by/4.0/, as recorded in the arXiv licence metadata for this exact version and shown on the abstract page https://arxiv.org/abs/2606.22401v1. Full licence notice: \texttt{licenses/arxiv-2606.22401-CC-BY-4.0.txt}.\par\smallskip

\noindent\textit{Editorial scope.} Counted unit: the article's lemma R[X\_m,M] (source0.tex lines 154--160). The article's complete original proof of it is delivered with it (source0.tex lines 161--234, one \texttt{proof} environment of 74 lines). No mathematics is written, repaired or re-derived: the statement and the proof are the article's own lines, in the article's own notation, and no retained line is substituted or re-flowed.  No further source-local context is needed: every cross-reference of every retained line resolves inside the two delivered spans.  Nothing was omitted from any retained span, so the package carries no omission record.  No retained line contains a citation, so the wrapper carries no bibliography and no bibliography entry.  No retained line contains a figure environment, an image-inclusion command or drawing code, so no image is delivered and the package declares no asset dependency.  Editorial formatting only, in addition: an AMS \texttt{amsart} preamble that restates the article's own declarations and macros used by the retained lines, namely the environments \texttt{lemma} on separate counters, together with the standard packages those lines need.  The retained environments print in this document as Lemma 1 in order of appearance, whereas the article prints the same items with the numbering of its own full document.  The complete native source of the article as distributed in the arXiv e-print for this exact version is bundled unchanged as \texttt{sources/203343/source0.tex}.
\begin{lemma} \label{LEM: zachowywanie długości przy przejściu z R^2 na R[X_m,M]}
Let $m<M$ be a pair of coprime positive integers and $X_{m,M}$ be the corresponding curve.  Let  $$\varphi: \R[x,y] \twoheadrightarrow \R[X_{m,M}]$$ be the natural projection. If the degree of $f$ is strictly smaller than $m$, then
    $$
    \ell_{2n, \R[x,y]} (f) = \ell_{2n,\R[X_{m,M}]}\big(\varphi(f)\big).
    $$
   
\end{lemma}

\begin{proof}
    Since $\varphi$ is a surjection, inequality $\ell_{2n, \R[x,y]} (f) \geq \ell_{2n,\R[X_{m,M}]}\big(\varphi(f)\big)$ holds. To finish the proof it is therefore enough to show that opposite inequality holds. We will restrict ourselves to the nontrivial case, that is $\ell_{2n,\R[X_{m,M}]}\big(\varphi(f)\big) < \infty$.

   
First of all, the coordinate ring of $X_{m,M}$ is easily seen to be isomorphic to the ring $\R[t^m,t^M]$. To stress out that we work in this particular ring, rather than in $\R[X_{m,M}]$, we denote by $\widetilde{\varphi}$ the epimorphism $\widetilde{\varphi}: \R[x,y] \rightarrow \R[t^m,t^M]$, and so, we will consider $\widetilde{\varphi}(f)$ instead of $\varphi(f)$.

    
% https://q.uiver.app/#q=WzAsMyxbMCwwLCJcXFJbeCx5XSJdLFsxLDAsIlxcUlt0Xm0sdF5NXSJdLFsxLDEsIlxcUltYX3ttLE19XSJdLFswLDEsIlxcd2lkZXRpbGRle1xcdmFycGhpfSJdLFswLDIsIlxcdmFycGhpIiwyXSxbMiwxLCJcXGlvdGEiLDJdXQ==
Let $f \in \R[x,y]$ be such that $\deg f <m$. Since $\deg\big(\widetilde{\varphi}(x)\big), \deg\big(\widetilde{\varphi}(y)\big) \leq M$, we obtain the bound 
    $$
    \deg\big(
    \widetilde{\varphi}(f)
    \big) \leq M \deg(f) < m\cdot M.
    $$
    Let us denote $\ell := \ell_{2n,\R[t^m,t^M]}
    \big(
    \widetilde{\varphi}(f)
    \big)$ and let $g_1,\ldots,g_\ell \in \R[t^m,t^M]$ be such that
    $$
     \widetilde{\varphi}(f) = \sum_{i=1}^\ell g_i^{2n}.
    $$
    We can immediately observe that for any $i =1,2,\dots,\ell$ inequality $2n \cdot \deg(g_i) \leq \deg\big(
    \widetilde{\varphi}(f)
    \big) < m M$ holds, hence $\deg(g_i) < \frac{mM}{2n}$. 
    
    To finish the proof we will define a $\R$-linear map $\psi: \R[t^m,t^M] \rightarrow \R[x,y]$ which satisfies $\widetilde{\varphi} \circ \psi = \mathrm{id}_{\R[t^m,t^M]}$. If $t^k \in \R[t^m,t^M]$, then there exists $a_k,b_k \in \Z_{\geq 0}$ satisfying $k = a_k m + b_k M$. Let us therefore define $\psi(t^k) := x^{a_k} y^{b_k}$ and extend linearly. By construction,  $$\widetilde{\varphi} \circ \psi(t^k) = t^{a_k m + b_k M} = t^k$$ hence indeed equality $\widetilde{\varphi} \circ \psi = \mathrm{id}_{\R[t^m,t^M]}$ holds. Additionally, we have the following chain of inequalities
    $$
    \deg \psi(t^k) = a_k + b_k = \frac{1}{m} \big( a_k m + b_k m\big) \leq \frac{1}{m} \big(a_km +  b_kM\big) = \frac{k}{m} = \frac{1}{m} \cdot \deg(t^k).
    $$
    As a consequence, the inequality $\deg \psi(h) \leq \frac{1}{m} \cdot \deg(h)$ holds for any $h \in \R[t^m,t^M]$.

    
    Now let $\mathfrak{s} := \sum_{i=1}^{\ell} {\psi(g_i)}^{2n}$ and observe that the following inequalities holds
    $$
    \deg(\mathfrak{s}) \leq \max_{i =1,\dots,\ell} \Big\{ 
    2n \cdot \deg \psi(g_i) 
    \Big\} 
    \leq
    \max_{i =1,\dots,\ell} \bigg\{ 
    \frac{2n}{m} \cdot \deg(g_i) 
    \bigg\}
    <
    \max_{i =1,\dots,\ell} \bigg\{ 
    \frac{2n}{m} \cdot \frac{mM}{2n} 
    \bigg\} = M.
    $$
    Our aim now, is to show that $\mathfrak{s} = f$. To accomplish this, let $\varepsilon := f -\mathfrak{s}$. We have
    \begin{multline*}
    \widetilde{\varphi} (\varepsilon) = \widetilde{\varphi}(f) - \widetilde{\varphi}(\mathfrak{s}) = 
    \sum_{i=1}^\ell g_i^{2n} - 
    \widetilde{\varphi}
    \Bigg(
    \sum_{i=1}^{\ell} {\psi(g_i)}^{2n}
    \Bigg)
    =
    \\
    =
    \sum_{i=1}^\ell g_i^{2n} - 
    \sum_{i=1}^{\ell} 
    {
    \Big(
    (\widetilde{\varphi} \circ \psi)(g_i)
    \Big)}^{2n} = 
    \sum_{i=1}^\ell g_i^{2n} - \sum_{i=1}^\ell g_i^{2n} = 0,
    \end{multline*}
    hence $\varepsilon \in \ker (\widetilde{\varphi})$. However, it is a straightforward calculation that $\ker (\widetilde{\varphi}) = (x^M - y^m)$, hence $\deg(\varepsilon) \geq \deg(x^M - y^m) = M$ or $\varepsilon = 0$. Since $\deg(f) < M$ and $\deg(\mathfrak{s}) < M$ we obtain $\deg(\varepsilon) \leq \max \big\{ \deg(f), \deg(\mathfrak{s}) \big\} < M$ which implies $\varepsilon = 0$. This finishes the proof, since we obtain
    $$
    f = \mathfrak{s} = \sum_{i=1}^{\ell} {\psi(g_i)}^{2n},
    $$
    which shows that $\ell_{2n,\R[x,y]}(f) \leq \ell = \ell_{2n,\R[X_{m,M}]}
    \big(
    \widetilde{\varphi}(f)
    \big)$ as desired.
\end{proof}

\end{document}
